% X 替换为真实问题的中文序号；编号由统一生成器决定。
\section{问题X求解与分析}

\subsection{问题的描述与分析}

说明输入、输出、目标、约束和问题依赖。

% 先描述与分析并引用，再在本节插入合格的主要图文综合流程图。
% 该图必须在下一二级标题之前实际排入 PDF，不能只有源码顺序正确。
% 使用 comprehensive-graphic-text-flowchart，遵守结构图工作流，不能用文字框占位图交付。
% 本问另需 4--6 张不重复的独立科学图；主要综合图及内部面板不抵扣额度。

\subsection{模型建立与求解}

给出假设依据、变量、公式、算法与参数来源。

\subsection{结果分析与验证}

引用真实结构化结果，报告独立验证、误差和适用边界。

% 后续小节按科学内容命名，必要时展开真实运行的基线比较、独立互验、
% 误差/约束、灵敏度、鲁棒性、场景、复杂度、收敛或理论边界。
% 所有数字可追溯到结构化结果。每个 Figure 围绕独立科学问题组织互补 Panel；
% 正式 caption 为实际 PDF 中含编号的一行短标题，长说明放正文或独立图注。
